\chapter{Anticonceptie}\label{ch:g8:contraception}

De vorige hoofdstukken brachten een mechanisme van opmerkelijke kracht in
kaart: twee cellen treffen elkaar, en veertig weken later is er een mens.
Mensen zijn de enige \omterm{def:g5:biodiversity:species}{soort} die haar eigen mechanisme begrijpt --- en
begrijpen brengt een keuze mee die de kikker nooit voorgelegd krijgt: of
en wanneer je het laat draaien. \emph{\omterm{prop:g8:contraception:principle}{Anticonceptie}} is de toegepaste
wetenschap van die keuze, en zij werkt, in elke methode, door een stap te
onderbreken die je nu kunt benoemen.

\section{Het principe: benoem de schakel, blokkeer de schakel}

\begin{proposition}[Anticonceptie past op het mechanisme]\label{prop:g8:contraception:principle}
Een \omterm{def:g5:human-reproduction:pregnancy}{zwangerschap} vraagt een ononderbroken keten
(\cref{prop:g8:sexual-reproduction:core},
\cref{prop:g8:from-egg-to-birth:firstweek}): \omterm{def:g8:sexual-reproduction:gametes}{gameten} gemaakt --- sperma
afgezet --- de zwemtocht naar de \omterm{def:g8:reproductive-systems:female}{eileider} --- een \omterm{def:g4:animal-reproduction:sexual}{eicel} losgelaten en
ontmoet --- \omterm{def:g5:flower-to-fruit:steps}{bevruchting} --- innesteling. \emph{Anticonceptie}\index{anticonceptie}
is elke bewuste onderbreking van die keten. Elke methode op de plank van
de apotheek is een aanval op één benoemde schakel, en ze op hun schakel
sorteren is de hele theorie:
\begin{itemize}
  \item \emph{barrière}methoden verhinderen dat de \omterm{def:g8:sexual-reproduction:gametes}{gameten} elkaar treffen;
  \item \emph{hormonale} methoden verhinderen dat er überhaupt een \omterm{def:g4:animal-reproduction:sexual}{eicel}
        wordt losgelaten;
  \item andere methoden laten de ontmoeting of de innesteling mislukken.
\end{itemize}
\end{proposition}
\admitted

\section{De belangrijkste methoden, op hun schakel}

\begin{definition}[Barrièremethoden: het condoom]\label{def:g8:contraception:condom}
Het \emph{condoom}\index{condoom} --- een dun omhulsel dat over de \omterm{def:g8:reproductive-systems:male}{penis}
gedragen wordt (een vrouwelijke versie bekleedt de vagina) --- blokkeert de
keten bij de stap van het afzetten: het sperma komt nooit in het lichaam
van de partner. Overal verkrijgbaar, en alleen nodig wanneer het nodig is
--- en uniek op de plank vanwege een tweede dienst: het is de enige
methode die ook de doorgang van de \emph{infecties} blokkeert die bij seks
overgedragen worden, het onderwerp van de wereld van
\cref{ch:g9:microbes-and-infection}. Die dubbele dienst is de reden dat het
condoom in beeld blijft, ook als er al een andere methode gebruikt wordt.
\end{definition}

\begin{definition}[Hormonale methoden: de pil]\label{def:g8:contraception:pill}
De \emph{anticonceptiepil}\index{anticonceptiepil} gebruikt de eigen
machinerie uit \cref{def:g8:puberty:hormones}: dagelijks ingenomen houden
haar \omterm{def:g8:puberty:hormones}{hormonen} de cyclus van de \omterm{def:g8:reproductive-systems:female}{eierstokken} in de opgeschorte toestand
(gebeurtenis 4 uit \cref{prop:g8:reproductive-systems:cycle}, maar zonder
enige \omterm{def:g5:human-reproduction:pregnancy}{zwangerschap}) --- \emph{geen \omterm{prop:g8:reproductive-systems:cycle}{eisprong}, geen \omterm{def:g4:animal-reproduction:sexual}{eicel}, niets om te
bevruchten}. Voorgeschreven door een arts, die de methode bij de persoon
zoekt; zeer betrouwbaar bij regelmatig innemen --- en zonder ook maar iets
tegen infecties te bieden, vandaar het vaste partnerschap met het \omterm{def:g8:contraception:condom}{condoom}.
\end{definition}

\begin{example}[De plank, op schakel gesorteerd]\label{ex:g8:contraception:shelf}
De andere gangbare methoden, op de aangevallen schakel gearchiveerd:
hormonale implantaten en pleisters --- het principe van de pil met een
beter geheugen (maanden of jaren, geen dagelijkse inname); de kleine
apparaatjes die een arts in de \omterm{prop:g5:human-reproduction:plan}{baarmoeder} plaatst --- die de ontmoeting en
de innesteling jarenlang laten mislukken; en de \emph{morning-afterpil},
een hormoondosis achteraf binnen enkele dagen die vooral werkt door een
\omterm{prop:g8:reproductive-systems:cycle}{eisprong} die nog niet gebeurd is uit te stellen --- een noodrem, bij de
apotheek verkrijgbaar, en geen vervanging van een methode die je
\emph{vooraf} gebruikt. Wat geen van deze doen: iets tegen infecties.
\end{example}

\begin{remark}[Wat niet werkt]\label{rem:g8:contraception:notwork}
De methoden van het schoolplein, afgemeten aan de keten: ``veilige dagen
tellen'' loopt stuk op de variabiliteit uit
\cref{rem:g8:reproductive-systems:living} --- de datum van de \omterm{prop:g8:reproductive-systems:cycle}{eisprong} is
vooraf niet te weten, en \omterm{def:g8:reproductive-systems:male}{zaadcellen} overleven er dagen op wachtend;
``op tijd terugtrekken'' loopt stuk op het feit dat \omterm{def:g8:reproductive-systems:male}{zaadcellen} al vóór de
finale onderweg zijn. Allebei falen ze precies waar de theorie van de
benoemde schakel het voorspelt: ze blokkeren geen enkele schakel
betrouwbaar. Een verhaal is geen methode.
\end{remark}

\section{Keuze, verantwoordelijkheid en de hulpkaart}

\begin{proposition}[De keuze is echt en gedeeld]\label{prop:g8:contraception:choice}
Feiten die het mechanisme oplegt:
\begin{itemize}
  \item een \omterm{def:g5:human-reproduction:pregnancy}{zwangerschap} heeft één keer nodig, niet de duur van een
        relatie: de keten telt geen gelegenheden;
  \item het is een \emph{gedeelde} keten: de cellen van allebei de
        partners, de verantwoordelijkheid van allebei --- het \omterm{def:g8:contraception:condom}{condoom} en
        de pil zijn niet elk de zaak van één geslacht, maar een gesprek;
  \item en geen enkele methode behalve onthouding is volmaakt, al komen
        de goede dicht in de buurt als ze precies volgens ontwerp gebruikt
        worden --- ``goed gebruikt'' is waar de meeste mislukkingen in de
        praktijk zitten.
\end{itemize}
\end{proposition}
\admitted

\begin{example}[De hulpkaart]\label{ex:g8:contraception:help}
De adressen uit \cref{ex:g8:puberty:questions}, uitgebreid:
schoolverpleegkundigen en huisartsen beantwoorden vragen over
anticonceptie vertrouwelijk --- ook die van minderjarigen; in elke regio
bestaan er centra voor gezinsplanning die daar juist voor zijn, gratis en
zonder oordeel; en apothekers bewaken de klok van de morning-afterpil. Het
systeem is zo gebouwd dat niemand de keten op verhalen hoeft te navigeren
--- de hulpkaart is een onderdeel van de methode.
\end{example}

\begin{method}[Elke methode beoordelen waarover je hoort]\label{met:g8:contraception:evaluate}
Vraag bij elke beweerde methode op volgorde:
\begin{enumerate}
  \item \emph{welke schakel} van de keten blokkeert ze --- kun je die op
        de kaart benoemen?
  \item \emph{hoe betrouwbaar}, zoals echte mensen haar gebruiken?
  \item beschermt ze ook tegen \emph{infecties}, of moet het \omterm{def:g8:contraception:condom}{condoom} haar
        vergezellen?
  \item wat vraagt ze --- dagelijks geheugen, een recept, een plaatsing
        --- en past dat bij de gebruiker?
\end{enumerate}
Geen te benoemen schakel, geen methode
(\cref{rem:g8:contraception:notwork} zakt al bij vraag 1).
\end{method}

\section{Oefeningen}

\begin{exercise}[$\star$]\label{exo:g8:contraception:1}
Beschrijf de keten die een \omterm{def:g5:human-reproduction:pregnancy}{zwangerschap} vraagt, en omschrijf anticonceptie
daartegenover.
\end{exercise}

\begin{exercise}[$\star$]\label{exo:g8:contraception:2}
Welke schakel blokkeert het \omterm{def:g8:contraception:condom}{condoom}? En de pil?
\end{exercise}

\begin{exercise}[$\star$]\label{exo:g8:contraception:3}
Welke tweede dienst levert alleen het \omterm{def:g8:contraception:condom}{condoom}?
\end{exercise}

\begin{exercise}[$\star$]\label{exo:g8:contraception:4}
Hoe werkt de morning-afterpil vooral, en wat is ze niet?
\end{exercise}

\begin{exercise}[$\star$]\label{exo:g8:contraception:5}
Noem drie adressen op de hulpkaart, met wat elk aanbiedt.
\end{exercise}

\begin{exercise}[$\star$]\label{exo:g8:contraception:6}
Waarom is ``veilige dagen tellen'' onbetrouwbaar? Twee redenen uit
\cref{rem:g8:contraception:notwork}.
\end{exercise}

\begin{exercise}[$\star\star$]\label{exo:g8:contraception:7}
Sorteer de plank uit \cref{ex:g8:contraception:shelf} op aangevallen
schakel, met \omterm{def:g8:contraception:condom}{condoom} en pil erbij.
\end{exercise}

\begin{exercise}[$\star\star$]\label{exo:g8:contraception:8}
Waarom zeggen artsen tegen jonge stellen zo vaak ``pil \emph{én}
\omterm{def:g8:contraception:condom}{condoom}''? Twee beschermingen, elk in één zin.
\end{exercise}

\begin{exercise}[$\star\star$]\label{exo:g8:contraception:9}
Leg met \cref{prop:g8:reproductive-systems:cycle} uit in welke toestand de
pil de cyclus houdt, en waarom dat een \omterm{def:g5:human-reproduction:pregnancy}{zwangerschap} voorkomt.
\end{exercise}

\begin{exercise}[$\star\star$]\label{exo:g8:contraception:10}
``De cellen van allebei de partners, de zaak van allebei de partners.''
Verdedig die zin tegenover de gewoonte om de anticonceptie aan één kant
van het stel over te laten.
\end{exercise}

\begin{exercise}[$\star\star$]\label{exo:g8:contraception:11}
Laat \cref{met:g8:contraception:evaluate} lopen over: (a) het hormonale
implantaat; (b) ``op tijd terugtrekken''.
\end{exercise}

\begin{exercise}[$\star\star\star$]\label{exo:g8:contraception:12}
``\omterm{prop:g8:contraception:principle}{Anticonceptie} is toegepaste fysiologie: elke methode op de plank is een
hoofdstuk van dit boek, gericht afgevuurd.'' Onderbouw dat met drie
methoden en hun drie hoofdstukken.
\end{exercise}

\section{Probleem: Het consult bij de gezinsplanning}

\begin{problem}[{Weekendprobleem --- een middag vol vragen bij het centrum
voor gezinsplanning}]\label{pb:g8:contraception:1}
De (verzonnen) middag van een consulent gezinsplanning: vier bezoekers,
vier situaties, en één ketenkaart aan de muur.

\medskip\noindent\textbf{Deel I --- De kaart aan de
muur.}
\begin{enumerate}
  \item Teken de wandkaart: de keten van \omterm{def:g8:sexual-reproduction:gametes}{gameten} tot innesteling, met de
        drie aanvalspunten uit
        \cref{prop:g8:contraception:principle} erop gemarkeerd.
  \item Bezoeker 1 vraagt: ``hoe kan één pil per dag iets tegenhouden?''
        Antwoord met de machinerie: wiens taal spreekt de pil, en tegen
        wie?
  \item Bezoeker 1 vraagt door: ``en als ik er een paar vergeet?'' Zoek
        het gevaar op de kaart --- wat hervat er als het hormoonsignaal
        wegvalt?
  \item Waarom overleeft de aanbeveling van het \omterm{def:g8:contraception:condom}{condoom} elke andere
        methodekeuze? (De dubbele dienst.)
\end{enumerate}

\medskip\noindent\textbf{Deel II --- De situaties.}
\begin{enumerate}[resume]
  \item Bezoeker 2, zeventien en ongerust, dacht dat een app met veilige
        dagen andere methoden overbodig maakte. Corrigeer die gedachte
        vriendelijk: welke twee biologische feiten verslaan de kalender?
  \item Bezoeker 3 wil iets ``zonder dagelijks geheugen''. Bied de twee
        kandidaten van de plank aan, met hun schakels.
  \item Bezoeker 4 komt binnen een dag na een gescheurd \omterm{def:g8:contraception:condom}{condoom}. Noem de
        mogelijkheid, haar klok, haar belangrijkste werking --- en de
        vaste methode die de consulent voor daarna zal bespreken.
  \item Alle vier vertrekken met dezelfde slotzin over infecties. Schrijf
        die op.
\end{enumerate}

\medskip\noindent\textbf{Deel III --- De nabespreking.}
\begin{enumerate}[resume]
  \item De stagiair vraagt waarom het centrum de keten onderwijst in
        plaats van gewoon methoden uit te delen. Antwoord met de eerste
        vraag van \cref{met:g8:contraception:evaluate}.
  \item Sorteer de methoden van die middag naar wat elk van haar gebruiker
        vraagt (vraag 4 van de methode).
  \item De stagiair vraagt zich af waarom bezoeken van minderjarigen
        standaard vertrouwelijk zijn. Antwoord met de logica van
        \cref{ex:g8:contraception:help}.
  \item Sluit de middag af in één zin: wat het centrum werkelijk
        verstrekt, naast doosjes.
\end{enumerate}
\end{problem}
